% !TeX root = ../main.tex
% !TeX encoding = UTF-8
% !TeX spellcheck = none

% 原文此节部分引文编号与文末条目不对应，已保留；见 conversion-notes.md。
\section{相关工作}\label{sec:related-work}

\subsection{时空建模与频率解耦}\label{sec:frequency-related}

交通流预测的核心任务是在复杂空间依赖和动态时间变化之间建立有效映射。基于图神经网络的方法通过显式建模道路节点之间的关联关系，成为交通预测的重要研究方向。DCRNN 将交通状态传播建模为有向图上的扩散过程，并结合扩散卷积与循环神经网络捕获空间传播和时间演化规律 \cite{ref01}。STGCN 进一步将图卷积与时间卷积结合，实现高效的时空特征提取 \cite{ref02}。Graph WaveNet 通过自适应邻接矩阵学习潜在空间关系，并利用扩张时间卷积扩大时间感受野，从而突破预定义拓扑限制 \cite{ref03}。这些研究建立了交通预测中空间关系建模与时间序列编码的基本范式。

随后，研究者进一步从动态空间关系和多尺度交互角度增强时空建模能力。ASTGCN 和 GMAN 分别利用空间---时间注意力机制动态刻画节点间相关性，使模型能够关注不同时间和空间位置的重要信息 \cite{ref04,ref05}。MTGNN 和 AGCRN 通过数据驱动方式学习节点间关系，降低了对人工构建邻接矩阵的依赖 \cite{ref06,ref07}。STSGCN 和 STFGNN 则从时空同步建模以及多尺度时空融合角度进一步提升复杂交通模式的表达能力 \cite{ref08,ref09}。这些方法有效增强了空间依赖学习能力，但其核心仍主要围绕统一时空表示学习展开。

近年来，频率解耦方法进一步关注交通序列中不同时间动态的差异性。STWave利用小波分解将交通信号划分为不同频率成分，并分别学习其时空依赖关系，从而减少不同动态模式之间的相互干扰 \cite{ref10}。TEDDCF \cite{ref18}则进一步区分趋势分量与事件分量，并为低频趋势和高频事件设计独立的时间与空间编码路径 \cite{ref11}。这类方法表明，交通流中的稳定演化规律与快速变化模式具有不同的统计特征，需要通过差异化建模方式进行刻画。

然而，现有时空建模与频率解耦方法主要关注如何获得更具表达能力的基础时空表示，而对于当交通状态偏离常规模式后，如何利用不同来源的补充信息进一步修正预测结果仍缺少深入研究。

\subsection{异质证据与专家建模}\label{sec:experts-related}

交通流演化通常同时包含周期规律、趋势变化、局部扰动以及节点特定行为，因此利用不同结构组件刻画异质动态成为重要研究方向。ST-ResNet 将城市人群流动划分为时间邻近性、周期性和趋势性因素，并采用多个残差分支分别进行建模 \cite{ref12}。该方法证明了不同时间模式具有不同预测作用。

进一步地，部分研究开始探索通过解耦或专家机制显式区分不同交通动态。D2STGNN 通过解耦扩散信号与内在信号，使不同来源的交通变化能够分别建模，并增强模型对复杂动态的表达能力 \cite{ref13}。TESTAM 则设计时间专家、静态图专家和动态图专家，并通过专家路由机制根据输入状态选择不同专家组合，以适应交通动态变化 \cite{ref14}。

上述研究表明，交通预测中的不同动态因素具有不同作用机制，采用具有专门归纳偏置的结构能够提升模型对复杂交通模式的建模能力。然而，现有专家化方法主要关注如何根据输入状态选择或融合不同专家，其目标通常是获得更加灵活的整体表示，而不是围绕“预测修正”任务组织证据。

对于交通状态偏离常规模式的情境，关键问题并非简单判断使用哪个专家，而是明确不同预测证据承担的修正角色。例如，周期模式用于提供稳定参照，局部变化用于捕获短时扰动，有向时滞关系用于描述跨区域历史关联，而节点自身演化用于保持局部连续性。这些证据具有不同时间语义和空间关系，需要采用不同的处理路径。

\subsection{动态依赖与历史消息选择}\label{sec:messages-related}

除了异质时间模式建模外，交通预测还需要考虑空间关系随时间变化以及跨区域历史信息的有效利用。注意力机制为动态依赖学习提供了一种数据驱动方式。GMAN 通过空间注意力和时间注意力捕获节点间动态相关性，并利用转换注意力建立历史状态与未来预测之间的关联 \cite{ref15}。PDFormer 则进一步考虑交通传播过程中的时延因素，通过延迟感知特征变换建模动态长程空间依赖 \cite{ref16}。这些研究说明，交通关系不仅由空间距离决定，还受到时间对齐和传播延迟影响。

与此同时，自适应图学习方法进一步探索动态空间结构。MTGNN、AGCRN 等方法通过学习节点关系替代固定邻接矩阵，使模型能够捕获潜在空间关联 \cite{ref06,ref07}。然而，这类方法通常关注节点关系的生成或调整，而较少进一步区分：在当前交通背景和预测时距条件下，某一历史消息是否真正具有预测价值。

事实上，跨区域历史信息的有效性具有明显的条件依赖性。同一邻居区域可能在不同历史时滞、不同时间尺度以及不同预测步下表现出不同参考价值。因此，跨区域建模不仅需要回答“哪些区域相关”，还需要回答“哪些历史模式应该被选择，以及选择程度如何随预测目标变化”。
